% Nouveau EGA V (cote 25): archive cote 25, site page 125 (PDF page 126).
% English translation of the checked French annotated master.
% AI-translated by GPT-6 Astra (OpenAI, ChatGPT Pro mode).
\scanpage{25/125}
\var{\inserted{taking into account that the hyp. on $f$ implies $f_*(\underline{O}_X)\simeq\underline{O}_S$, whence $f_*(\underline{O}_X^*)\simeq\underline{O}_S^*$}}{taking into account that the hyp. on $f$ implies $f_*(\underline{O}_X)\simeq\underline{O}_S$, whence $f_*(\underline{O}_X^*)\simeq\underline{O}_S^*$}{insertion}{ega25_25_125-v0001}
\mspage{\hand{28}}

by $\overset{\hand{\scriptscriptstyle\vee}}{P}$ , and associated with $\underline{L}$ (in the sense of \hand{4}.1.) correspond bijectively
to the\hand{*} \uline{surjective} homomorphisms $Q\hand{\longrightarrow}\overset{\hand{\scriptscriptstyle\vee}}{E}$\hand{,} modulo \var{\inserted{multiplication by}}{multiplication by}{insertion}{ega25_25_125-v0002} sections of $\underline{O}_S^{\hand{*}}$ .\hand{$\blacktriangledown$}
\var{\inserted{\var{\deleted{i.e. to the sections of the sheaf $\mathcal{H}om_S(Q,\overset{\hand{\scriptscriptstyle\vee}}{E})^*/\underline{O}_S^*$, where the \unsure{first} * \unsure{indicates}}}{}{deletion}{ega25_25_125-v0004}}}{\var{\deleted{i.e. to the sections of the sheaf $\mathcal{H}om_S(Q,\overset{\hand{\scriptscriptstyle\vee}}{E})^*/\underline{O}_S^*$, where the \unsure{first} * \unsure{indicates}}}{}{deletion}{ega25_25_125-v0004}}{insertion}{ega25_25_125-v0003}
\var{\margin{\unsure{For sheaves of sets} \illegible{2} \var{\deleted{\illegible{2}}}{}{deletion}{ega25_25_125-v0006} \unsure{specify}}}{\placedmargin{\unsure{For sheaves of sets} \illegible{2} \var{\deleted{\illegible{2}}}{}{deletion}{ega25_25_125-v0006} \unsure{specify}}}{marginal note placed}{ega25_25_125-v0005}
The existence of $Q$ is reduced by the patented procedure to the case of noetherian S,
and in that case is precisely III 7.7.6. (the hypothesis on the
fibers of X being moreover \unsure{unnecessary}).\tnote{25-125-1}{The French word ends in an uncertain e or è; the lexical reading is ``inutile''. The uncertainty remains marked in the translation.}
Since $\underline{E}$ is loc free of finite type,
giving a hom $\underline{E}\hand{\longrightarrow} f(\underline{L})$ is equivalent to giving a section of
$\underline{L}\otimes_{\underline{O}_S}\overset{\hand{\scriptscriptstyle\vee}}{\underline{E}}$ , hence to giving a homomorphism $Q\hand{\longrightarrow}\overset{\hand{\scriptscriptstyle\vee}}{\underline{E}}$ . It remains to express
that the condition considered in \hand{4}.1. is indeed satisfied, which (thanks to
the hypothesis made on the fibers of X/S) reduces to checking that
fiber by fiber the corresponding homomorphism $E(s)\hand{\longrightarrow}$ \var{\xout{Q(s)}}{}{typewriter strike-out}{ega25_25_125-v0007}\hand{$\vee$}
$H^0(X_s,L_s)\simeq\mathrm{Hom}_{k(s)}(Q(s),k(s))$ is injective, or again that $Q(s)\hand{\longrightarrow}\overset{\hand{\scriptscriptstyle\vee}}{E}(s)$ is surjective,
which by Nakayama also means that $Q$ \quad $E$ is surjective. The ``modulo
sections of $\underline{O}_S$ '' comes from the ``modulo isomorphism'' in \hand{4}.1.\hand{/}

One can \var{\xout{xxxxxxxxx}}{}{typewriter strike-out}{ega25_25_125-v0008}\var{\inserted[typed]{interpret}}{interpret}{insertion}{ega25_25_125-v0009} \hand{4}.1. by another route, using the fact
that $P(\underline{Q})$ represents the \var{\xout{xxxxxxxx}}{}{typewriter strike-out}{ega25_25_125-v0010} subfunctor of $\mathrm{Div}_{X/S}$ defined by
$\underline{L}$, by virtue of n$^{\circ}$2. Consequently, a linear system of divisors parametr
by $\overset{\hand{\scriptscriptstyle\vee}}{P}$ \var{\inserted[typed]{and associated with $\underline{L}$}}{and associated with $\underline{L}$}{insertion}{ega25_25_125-v0012} is interpreted as a morphism $\overset{\hand{\scriptscriptstyle\vee}}{P} = P(\overset{\hand{\scriptscriptstyle\vee}}{\underline{E}})\hand{\longrightarrow} P(Q)$; the
\uline{linear} character of the family of divisors defined by such a morphism
is then interpreted by the fact that this morphism is ``linear'', i.e.
precisely defined by a surjective morphism $Q\hand{\longrightarrow}\overset{\hand{\scriptscriptstyle\vee}}{\underline{E}}$ . \var{\margin{$\ulcorner$}}{\placedmargin{$\ulcorner$}}{marginal note placed}{ega25_25_125-v0013} One also sees
in\var{\xout{xx}}{}{typewriter strike-out}{ega25_25_125-v0014} this case that the morphism $\overset{\hand{\scriptscriptstyle\vee}}{P}\hand{\longrightarrow}\mathrm{Div}_{X/S}$ is a monomorphism \var{\deleted{\illegible{2}}}{}{deletion}{ega25_25_125-v0015}
(since $\overset{\hand{\scriptscriptstyle\vee}}{P}\hand{\longrightarrow} P(Q)\hand{\hookrightarrow}\mathrm{Div}_{X/S}$ is one), a \var{\inserted{more}}{more}{insertion}{ega25_25_125-v0016} general fact moreover, cf. the corollary
below. Let us \var{\inserted[typed]{then}}{then}{insertion}{ega25_25_125-v0017} agree to say that \unsure{two} linear families of divisors
of X/S, parametrized by the projective bundles $\overset{\hand{\scriptscriptstyle\vee}}{P},\overset{\hand{\scriptscriptstyle\vee}}{P}'$, are \uline{isomorphic}
when they are transformed into each other by an S-isomorphism $\overset{\hand{\scriptscriptstyle\vee}}{P}\hand{\longrightarrow}\overset{\hand{\scriptscriptstyle\vee}}{P}'$ ,
(which will moreover be unique thanks to the fact that there are monomorphisms
into $\mathrm{Div}_{X/S}$ ). \var{\deleted{We shall call a linear system of divisors in X/S}}{}{deletion}{ega25_25_125-v0018}
(\var{\deleted{without reference to a specified projective bundle parametrizing it) a}}{}{deletion}{ega25_25_125-v0019}
\tnote{25-125-2}{In the French word translated as ``two'', the second letter is uncertain between e and a. The clipped ``paramétr'' earlier in the paragraph is retained as the fragment ``parametr''. The source abbreviations ``loc'' and ``n°'' mean locally and number. The manuscript phrase ``procédé brévété'' is rendered literally as ``patented procedure''.}

\mspage{\hand{125}}
